\chapter{电子量子测量与可辨识性}\label{chap:electron-quantum-measurement}

约化电子态尚未给出探测器实际记录的概率。测量仪器把联合态送到可读结果，扫描参考操作和统计信息再决定哪些相位及参数能被实验区分。

\section{量子测量与统计信息}

已经求得探测面上的密度算符后，接下来计算探测器会记录什么。理想能量测量给出某条边带的概率，实际谱仪则把有限能量区间分配给同一个读数，还可能漏掉部分电子。我们先从熟悉的投影测量出发，建立这些读出过程的概率公式，再讨论数据能够确定哪些态分量，以及计数有限时参数能测多准。

\subsection{量子仪器与正算符值测度（POVM）}

回忆投影测量。若一组正交态 $|x\rangle$ 对应可区分的离散结果，投影算符为 $P_x=|x\rangle\langle x|$，结果概率为 $\operatorname{Tr}(P_x\rho_S)$。真实谱仪的一个读数可能来自几个不同能量，因而该结果需要由更一般的正算符 $\Pi_x^{(S)}$ 描述。这里“正”指对任意态 $|\psi\rangle$，都有 $\langle\psi|\Pi_x^{(S)}|\psi\rangle\ge0$。一组和为恒等算符的这种效应算符称为\emph{正算符值测度}（POVM）；它们保证各结果概率非负且总和为一。

如果还要知道读到 $x$ 后系统处于什么态，就需要规定状态更新。用线性映射 $\mathcal I_x$ 表示这个更新，但暂不把输出归一化：它的迹保留结果发生的概率。所有结果对应的一族映射称为\emph{量子仪器}。每个 $\mathcal I_x$ 必须完全正，即只对更大纠缠系统中的 $S$ 部分作此更新，联合态仍保持非负；它还必须迹不增，因为单个结果的概率不能超过一。把所有结果加起来则保迹。因此
\begin{equation}
\boxed{
p(x)=\operatorname{Tr}[\mathcal I_x(\rho_S)],
\qquad
\rho_{S|x}=\frac{\mathcal I_x(\rho_S)}{p(x)},
\qquad
\sum_x\mathcal I_x\ \text{保迹}.}
\label{eq:quantum_instrument_probability_definition}
\end{equation}
当 $p(x)>0$ 时，除以迹就得到归一化条件态 $\rho_{S|x}$；零概率结果不定义条件态。为了只计算概率，可以把映射从态移到观测算符上。定义\emph{对偶映射} $\mathcal I_x^*$ 满足
\[
\operatorname{Tr}[A\mathcal I_x(\rho)]
=\operatorname{Tr}[\mathcal I_x^*(A)\rho]
\]
对所考虑的任意 $A,\rho$ 成立。取 $A=\mathbb I_S$，就得到效应算符 $\Pi_x^{(S)}=\mathcal I_x^*(\mathbb I_S)$，以及
\begin{equation}
\boxed{
\Pi_x^{(S)}\succeq0,
\qquad
\sum_x\Pi_x^{(S)}=\mathbb I_S,
\qquad
p(x)=\operatorname{Tr}[\Pi_x^{(S)}\rho_S].}
\label{eq:abstract_instrument_to_povm}
\end{equation}
同一个 $\Pi_x^{(S)}$ 可以来自不同的 $\mathcal I_x$，所以相同统计并不唯一决定测量反作用。

\begin{theorem}[量子仪器诱导 POVM]
设 $\{\mathcal I_x\}$ 是有限维 CP、迹不增仪器，$\sum_x\mathcal I_x$ 保迹。则 $\Pi_x=\mathcal I_x^*(I)$ 为正效应，$\sum_x\Pi_x=I$，且 $p(x)=\Tr(\Pi_x\rho)$。同一组效应不唯一决定条件态反作用。
\end{theorem}
\begin{proof}
量子仪器的每个结果分支 $\mathcal I_x$ 是完全正、迹不增映射，而无条件映射
\[
\mathscr E=\sum_x\mathcal I_x
\]
保迹。定义对偶映射 $\mathcal I_x^*$ 使任意系统算符 $X$ 与迹类算符 $\rho$ 满足
\[
\Tr[X\,\mathcal I_x(\rho)]
=\Tr[\mathcal I_x^*(X)\rho].
\]
取 $X=\mathbb I_S$，定义
\[
\Pi_x^{(S)}\equiv\mathcal I_x^*(\mathbb I_S).
\]
于是结果概率逐步变为
\begin{align*}
p(x)
&=\Tr[\mathcal I_x(\rho_S)]\\
&=\Tr[\mathbb I_S\mathcal I_x(\rho_S)]\\
&=\Tr[\mathcal I_x^*(\mathbb I_S)\rho_S]\\
&=\Tr[\Pi_x^{(S)}\rho_S].
\end{align*}

再验证 $\Pi_x^{(S)}$ 是正算符。任取 $|\psi\rangle$，令
$\rho_\psi=|\psi\rangle\langle\psi|\succeq0$。完全正性蕴含正性，因此
$\mathcal I_x(\rho_\psi)\succeq0$，从而
\begin{align*}
\langle\psi|\Pi_x^{(S)}|\psi\rangle
&=\Tr[\Pi_x^{(S)}\rho_\psi]\\
&=\Tr[\mathcal I_x(\rho_\psi)]\\
&\ge0.
\end{align*}
该式对任意 $|\psi\rangle$ 成立，所以
\[
\Pi_x^{(S)}\succeq0.
\]
迹不增还给出 $\Pi_x^{(S)}\preceq \mathbb I_S$。因为对任意密度算符
\[
0\le\Tr[\mathcal I_x(\rho)]\le\Tr\rho,
\]
所以
\[
0\le\Tr[(\mathbb I_S-\Pi_x^{(S)})\rho]
\]
对任意 $\rho\succeq0$ 成立，故 $\mathbb I_S-\Pi_x^{(S)}\succeq0$。

最后利用无条件映射保迹：
\begin{align*}
\Tr\rho
&=\Tr[\mathscr E(\rho)]\\
&=\sum_x\Tr[\mathcal I_x(\rho)]\\
&=\Tr\!\left[\left(\sum_x\Pi_x^{(S)}\right)\rho\right].
\end{align*}
因此
\[
\Tr\!\left[\left(\mathbb I_S-\sum_x\Pi_x^{(S)}\right)\rho\right]=0
\]
对任意迹类 $\rho$ 成立，只能有
\[
\sum_x\Pi_x^{(S)}=\mathbb I_S.
\]
于是 $\{\Pi_x^{(S)}\}$ 满足正性、完备性和 Born 概率公式，正是正文\eqnrefs{eq:abstract_instrument_to_povm}。
\end{proof}


\begin{proposition}[完全正仪器、效应算符与条件态]
在有限工作空间中，若测量结果 $x$ 的更新由 Kraus 算符
$M_{x\alpha}$ 给出，且 $\sum_{x,\alpha}M_{x\alpha}^\dagger M_{x\alpha}=I$，则
\[
\mathcal I_x(\rho)=\sum_\alpha M_{x\alpha}\rho M_{x\alpha}^\dagger,
\qquad \Pi_x=\sum_\alpha M_{x\alpha}^\dagger M_{x\alpha}
\]
构成量子仪器和 POVM。每个分支完全正且迹不增，
$p_x=\Tr(\Pi_x\rho)$，所有分支之和保迹。
\end{proposition}
\begin{proof}
将任意参考系统 $R$ 加入，对任意联合正态 $X_{SR}$，
\[
(\mathcal I_x\otimes\mathcal I_R)(X_{SR})
=\sum_\alpha(M_{x\alpha}\otimes I_R)X_{SR}
(M_{x\alpha}^\dagger\otimes I_R)\succeq0.
\]
这直接证明完全正性。迹的循环性给
$\Tr\mathcal I_x(\rho)=\Tr(\Pi_x\rho)$。
完备关系给 $0\preceq\Pi_x\preceq I$ 与 $\sum_x\Pi_x=I$，因此分支迹不增且总迹保持。
例如两个仪器可取 $M_x=\sqrt{\Pi_x}$ 或 $M_x'=U_x\sqrt{\Pi_x}$，
它们的效应相同，但结果相关的幺正 $U_x$ 会改变条件态。
所以能谱概率不能独自确定后续测量的统计。
\end{proof}



抽象的 $\mathcal I_x$ 并不是额外假设出来的“测量公设”。实验上更具体的图像是：待测系统先与一个可控测量仪发生普通量子力学的联合幺正演化，然后只读取测量仪的一部分信息。忽略测量仪对应偏迹，选择某个测量仪结果对应条件化；Kraus 算符和 POVM 都从这两个操作中产生。物理实验常通过“系统 $S$ + 测量仪 $M$ + 联合幺正 + 测量仪读出”实现量子仪器。假设制备时系统与测量仪因子化，$\rho_{SM}^{(i)}=\rho_S\otimes\rho_M$，并把
$\rho_M=\sum_\nu p_\nu|\nu\rangle\langle\nu|$。若联合演化为 $U$，而实验完全忽略测量仪，系统先由联合态的偏迹定义无条件演化
\begin{equation}
\boxed{
\mathscr E(\rho_S)=\operatorname{Tr}_M\!\left[U(\rho_S\otimes\rho_M)U^\dagger\right].}
\label{eq:unitary_dilation_partial_trace_definition}
\end{equation}
对测量仪初态与末态基展开以后，它化为完全正保迹（CPTP）的 Kraus 形式
\begin{equation}
\boxed{
\mathscr E(\rho_S)
=\sum_{\mu,\nu}K_{\mu\nu}\rho_SK_{\mu\nu}^\dagger,
\qquad
K_{\mu\nu}=\sqrt{p_\nu}\langle\mu|U|\nu\rangle,
\qquad
\sum_{\mu,\nu}K_{\mu\nu}^\dagger K_{\mu\nu}=\mathbb I_S.}
\label{eq:general_unitary_dilation_kraus_operator}
\end{equation}
初始因子化是得到统一系统映射的物理条件；若 $S$ 与 $M$ 预先相关，一般不存在作用于任意 $\rho_S$ 的同一个 CPTP 映射，必须保留初始联合态，或改用过程张量/赋值映射描述。

\begin{theorem}[因子化初态的幺正扩张与 Kraus 算符]
设系统与测量仪初态为 $\rho_S\otimes\rho_M$，$\rho_M=\sum_\nu p_\nu|\nu\rangle\langle\nu|$，联合演化 $U$ 幺正。则偏迹映射为 CPTP，其 Kraus 算符及完备关系为\eqnrefs{eq:general_unitary_dilation_kraus_operator}。在可分无限维空间中，对密度算符的求和按迹范数极限理解。
\end{theorem}
\begin{proof}
设测量仪初态作谱分解
\[
\rho_M=\sum_\nu p_\nu|\nu\rangle\langle\nu|,
\qquad
p_\nu\ge0,
\qquad
\sum_\nu p_\nu=1,
\]
并在测量仪末态空间选取任意完备正交基 $\{|\mu\rangle\}$。偏迹的基展开为
\[
\Tr_M X=\sum_\mu\langle\mu|X|\mu\rangle_M.
\]
因此正文\eqnrefs{eq:unitary_dilation_partial_trace_definition} 逐步化为
\begin{align*}
\mathscr E(\rho_S)
&=\sum_\mu
\langle\mu|U(\rho_S\otimes\rho_M)U^\dagger|\mu\rangle\\
&=\sum_{\mu,\nu}p_\nu
\langle\mu|U(\rho_S\otimes|\nu\rangle\langle\nu|)U^\dagger|\mu\rangle.
\end{align*}
系统算符 $\rho_S$ 不作用于测量仪空间，因此
\begin{align*}
&\langle\mu|U(\rho_S\otimes|\nu\rangle\langle\nu|)U^\dagger|\mu\rangle\\
&\qquad=
\langle\mu|U|\nu\rangle\,
\rho_S\,
\langle\nu|U^\dagger|\mu\rangle.
\end{align*}
定义只作用于系统 Hilbert 空间的 Kraus 算符
\[
K_{\mu\nu}
\equiv\sqrt{p_\nu}\,\langle\mu|U|\nu\rangle,
\]
便有
\[
\mathscr E(\rho_S)
=\sum_{\mu\nu}K_{\mu\nu}\rho_SK_{\mu\nu}^\dagger.
\]

完全正性不能只检验系统自身的正性，还要允许系统与任意参考空间 $R$ 纠缠。对任意
$X_{SR}\succeq0$，扩张映射为
\[
(\mathscr E\otimes\mathbb I_R)(X_{SR})
=\sum_{\mu\nu}
(K_{\mu\nu}\otimes I_R)X_{SR}
(K_{\mu\nu}^\dagger\otimes I_R).
\]
每一项都具有 $YXY^\dagger$ 形式，故为正算符；正算符之和仍为正。因此 $\mathscr E$ 对任意辅助系统都保持正性，即完全正。

再验证迹保持。先计算 Kraus 完备关系：
\begin{align*}
\sum_{\mu\nu}K_{\mu\nu}^\dagger K_{\mu\nu}
&=\sum_{\mu\nu}p_\nu
\langle\nu|U^\dagger|\mu\rangle
\langle\mu|U|\nu\rangle\\
&=\sum_\nu p_\nu
\langle\nu|U^\dagger
\left(\sum_\mu|\mu\rangle\langle\mu|\right)
U|\nu\rangle\\
&=\sum_\nu p_\nu\langle\nu|U^\dagger U|\nu\rangle\\
&=\sum_\nu p_\nu \mathbb I_S\\
&=\mathbb I_S.
\end{align*}
于是
\begin{align*}
\Tr\mathscr E(\rho_S)
&=\sum_{\mu\nu}\Tr(K_{\mu\nu}\rho_SK_{\mu\nu}^\dagger)\\
&=\Tr\!\left[\left(\sum_{\mu\nu}K_{\mu\nu}^\dagger K_{\mu\nu}\right)\rho_S\right]\\
&=\Tr\rho_S.
\end{align*}
所以该映射同时完全正且保迹，得到正文\eqnrefs{eq:general_unitary_dilation_kraus_operator}。
\end{proof}


读出测量仪时，令结果 $x$ 对应测量仪侧效应 $\Pi_x^{(M)}\succeq0$ 且 $\sum_x\Pi_x^{(M)}=\mathbb I_M$。对测量仪取偏迹后，系统的未归一化条件态为
\begin{equation}
\boxed{
\mathcal I_x(\rho_S)
=\operatorname{Tr}_M\!\left[
(\mathbb I_S\otimes [\Pi_x^{(M)}]^{1/2})
U(\rho_S\otimes\rho_M)U^\dagger
(\mathbb I_S\otimes [\Pi_x^{(M)}]^{1/2})
\right].}
\label{eq:general_indirect_measurement_instrument}
\end{equation}
由于最终对测量仪取偏迹，两个 $[\Pi_x^{(M)}]^{1/2}$ 也可以利用偏迹的循环性合成一个 $\Pi_x^{(M)}$。对 $\Pi_x^{(M)}$ 作谱分解后，\eqnrefs{eq:general_indirect_measurement_instrument} 化为 Kraus 求和，因此每个分支完全正。系统对应的 POVM 效应算符为
\[
\Pi_x^{(S)}=\mathcal I_x^*(\mathbb I_S),
\]
而不是测量仪上的 $\Pi_x^{(M)}$ 本身。

\begin{proposition}[测量仪效应诱导的条件 CP 分支]
在上述因子化扩张中，令测量仪效应 $\Pi_x^{(M)}\succeq0$ 且和为单位。则\eqnrefs{eq:general_indirect_measurement_instrument}定义 CP、迹不增分支，所有分支之和等于无条件偏迹信道；系统效应为 $\mathcal I_x^*(I_S)$。最终测量仪被丢弃时，系统分支只依赖测量仪效应。
\end{proposition}
\begin{proof}
联合幺正结束以后，系统与测量仪的联合态为
\[
\rho'_{SM}=U(\rho_S\otimes\rho_M)U^\dagger.
\]
测量仪结果 $x$ 由效应 $\Pi_x^{(M)}\succeq0$ 描述，且
\[
\sum_x\Pi_x^{(M)}=\mathbb I_M.
\]
取正平方根
\[
M_x=[\Pi_x^{(M)}]^{1/2},
\qquad
M_x^\dagger M_x=\Pi_x^{(M)}.
\]
在测量仪上实施该结果对应的 L\"uders 操作后，联合未归一化分支为
\[
\widetilde\rho'_{SM|x}
=(\mathbb I_S\otimes M_x)\rho'_{SM}
(\mathbb I_S\otimes M_x^\dagger).
\]
只保留系统意味着对 $M$ 取偏迹，因此定义
\[
\mathcal I_x(\rho_S)=\Tr_M\widetilde\rho'_{SM|x}.
\]
把 $\rho'_{SM}$ 代回：
\begin{align*}
\mathcal I_x(\rho_S)
&=\Tr_M\!\left[
(\mathbb I_S\otimes M_x)
U(\rho_S\otimes\rho_M)U^\dagger
(\mathbb I_S\otimes M_x^\dagger)
\right]\\
&=\Tr_M\!\left[
(\mathbb I_S\otimes [\Pi_x^{(M)}]^{1/2})
U(\rho_S\otimes\rho_M)U^\dagger
(\mathbb I_S\otimes [\Pi_x^{(M)}]^{1/2})
\right],
\end{align*}
得到正文\eqnrefs{eq:general_indirect_measurement_instrument}。

该分支的迹就是结果概率。利用偏迹与全迹的关系以及循环性，
\begin{align*}
p(x)
&=\Tr_S\mathcal I_x(\rho_S)\\
&=\Tr_{SM}\!\left[
(\mathbb I_S\otimes M_x^\dagger M_x)
\rho'_{SM}
\right]\\
&=\Tr_{SM}\!\left[
(\mathbb I_S\otimes\Pi_x^{(M)})
U(\rho_S\otimes\rho_M)U^\dagger
\right].
\end{align*}
对所有结果求和时，
\begin{align*}
\sum_x\mathcal I_x(\rho_S)
&=\Tr_M\!\left[
\sum_x(\mathbb I_S\otimes M_x)
\rho'_{SM}
(\mathbb I_S\otimes M_x^\dagger)
\right],
\end{align*}
其迹满足
\begin{align*}
\Tr_S\sum_x\mathcal I_x(\rho_S)
&=\Tr_{SM}\!\left[
(\mathbb I_S\otimes\sum_x\Pi_x^{(M)})\rho'_{SM}
\right]\\
&=\Tr_{SM}\rho'_{SM}=1.
\end{align*}
完全正性可直接写成 Kraus 形式：
\[
 K_{x\mu\nu}=\sqrt{p_\nu}\langle\mu|(I_S\otimes M_x)U|\nu\rangle,
 \qquad \mathcal I_x(\rho_S)=\sum_{\mu\nu}K_{x\mu\nu}\rho_SK_{x\mu\nu}^\dagger.
\]
若测量仪使用另一组局域更新算符 $M_{x\alpha}$，只要
$\sum_\alpha M_{x\alpha}^\dagger M_{x\alpha}=\Pi_x^{(M)}$，对任意系统观测量 $O_S$ 都有
\[
 \Tr\!\left[O_S\sum_\alpha\Tr_M
 ((I\otimes M_{x\alpha})\rho'_{SM}(I\otimes M_{x\alpha}^\dagger))\right]
 =\Tr[(O_S\otimes\Pi_x^{(M)})\rho'_{SM}].
\]
右端不依赖局域更新的具体分解。因此测量仪被丢弃后的系统条件分支只依赖效应；
对结果求和便是 $\Tr_M\rho'_{SM}$，不仅是迹相同。
每个分支的效应不超过单位，所以迹不增，且所有分支之和就是原无条件偏迹信道。
\end{proof}


实验探测器通常还在理想量子结果之后加入有限分辨率和效率。对理想电子边带 $|\ell\rangle$，若能谱仪输出连续读数 $E_d$，定义归一化条件响应密度 $R(E_d|\mathcal E_\ell)$。当仪器只做经典能量模糊而不相干混合边带时，实际 POVM 密度为
\begin{equation}
\boxed{
\Pi_{\rm e}(E_d)
=\sum_\ell R(E_d|\mathcal E_\ell)|\ell\rangle\langle\ell|.}
\label{eq:electron_detector_response_povm}
\end{equation}
于是 $p(E_d)=\operatorname{Tr}[\Pi_{\rm e}(E_d)\rho_e]$。只有当 $R(E_d|\mathcal E_\ell)$ 还具有平移不变形式时，这个映射才退化为通常的卷积；能量依赖效率、角度分辨或相干混合都应直接进入更一般的 POVM，而不是修改相互作用哈密顿量。

\begin{exbox}[弱测量与弱值]
取有界无量纲系统可观测量 $\hat A$、居中的实高斯探针及固定的非正交预选、后选态。
求弱脉冲后的条件位移、成功概率，并说明近正交后选择怎样改变展开条件。

\textbf{解。}
\eqnrefs{eq:general_indirect_measurement_instrument} 已经给出一般“系统--探针--读出”测量。弱测量只是在其中选择一个特定的脉冲耦合并对耦合强度作受控展开。令系统的无量纲厄米可观测量为 $\hat A$，探针具有正则变量
\[
[\hat Q,\hat P]=\ii\hbar,
\]
并取脉冲相互作用
\begin{equation}
\hat H_{\rm int}(t)=\lambda(t)\hat A\otimes\hat P,
\qquad
g\equiv\int\dd t\,\lambda(t),
\qquad [g]={\rm m}.
\label{eq:weak_measurement_von_neumann_coupling}
\end{equation}
忽略脉冲期间的自由演化时，联合幺正为
\[
\hat U_w=\exp\!\left(-\frac{\ii g}{\hbar}\hat A\otimes\hat P\right).
\]
系统预选为 $|i\rangle$，探针初态为 $|\phi\rangle$，随后把系统后选择到 $|f\rangle$。当 $\langle f|i\rangle\neq0$ 时，定义无量纲弱值
\begin{equation}
\boxed{
A_w\equiv\frac{\langle f|\hat A|i\rangle}{\langle f|i\rangle}.}
\label{eq:weak_value_definition}
\end{equation}
若探针动量均方根宽度为 $\sigma_P$，一阶弱展开至少要求
\[
\epsilon_w^{(1)}
\equiv\frac{|g|\sigma_P}{\hbar}|A_w|\ll1,
\qquad
\epsilon_w^{(2)}
\equiv\frac{g^2\sigma_P^2}{2\hbar^2}
\left|\frac{\langle f|\hat A^2|i\rangle}{\langle f|i\rangle}\right|\ll1.
\]
第二个条件检查被 $\langle f|i\rangle^{-1}$ 放大的二阶矩，但这两个矩也可能因相消而同时很小。
对有界 $A$，完整一阶余项还可直接控制。令 $c_{fi}=\langle f|i\rangle$，由自伴算符的积分余项得
\[
 \left\|\left(U_w-I+\frac{\ii g}{\hbar}A\otimes P\right)|i,\phi\rangle\right\|
 \le\frac{g^2\|A\|^2}{2\hbar^2}\sqrt{\langle P^4\rangle_i}.
\]
高斯探针有 $\langle P^4\rangle_i=3\sigma_P^4$，因此一个直接检验归一化后选择态的充分条件为
\[
 \epsilon_{w,{\rm rem}}=
 \frac{\sqrt3\,g^2\|A\|^2\sigma_P^2}{2\hbar^2|c_{fi}|}\ll1,
 \qquad \epsilon_w^{(1)}\ll1.
\]
余项范数除以零阶成功振幅，而不是只比较未归一化联合态的误差。
若 $c_{fi}$ 也随 $g$ 趋零，应重新检验这两个量，不能沿用固定后选态下的 $O(g^2)$ 估计。对 $\langle Q\rangle_i=\langle P\rangle_i=0$、实且无 $Q$--$P$ 相关的最小不确定 高斯 探针，
$\sigma_P=\hbar/(2\sigma_Q)$，后选择后的平均读数为
\begin{equation}
\boxed{
\Delta\langle Q\rangle_f
=g\,\operatorname{Re}A_w+O(g^2),
\qquad
\Delta\langle P\rangle_f
=\frac{2g\sigma_P^2}{\hbar}\operatorname{Im}A_w+O(g^2)
=\frac{g\hbar}{2\sigma_Q^2}\operatorname{Im}A_w+O(g^2).}
\label{eq:weak_value_gaussian_meter_shifts}
\end{equation}
因此 $\operatorname{Re}A_w$ 与 $\operatorname{Im}A_w$ 分别控制共轭探针变量的一阶条件位移。$A_w$ 可以落在 $\hat A$ 的本征值区间之外，但它是带后选择条件的响应系数；后选择概率和高阶修正必须与位移同时计算。若 $\epsilon_w^{(1)}$ 或 $\epsilon_w^{(2)}$ 不小，应直接使用\eqnrefs{eq:general_indirect_measurement_instrument} 的完整量子仪器，而不使用弱值截断。

\medskip

后选择后的未归一化探针态由系统偏矩阵元给出：
\[
|\widetilde\phi_f\rangle
=\langle f|\hat U_w|i\rangle|\phi\rangle.
\]
对 $g$ 展开联合幺正，
\begin{align*}
\langle f|\hat U_w|i\rangle
&=\langle f|i\rangle
-\frac{\ii g}{\hbar}\langle f|\hat A|i\rangle\hat P
+O(g^2)\\
&=\langle f|i\rangle
\left(1-\frac{\ii g}{\hbar}A_w\hat P\right)+O(g^2).
\end{align*}
这个余项界来自
$e^{-\ii gX}-I+\ii gX=-\int_0^g\dd s\,(g-s)X^2e^{-\ii sX}$，
其中 $X=A\otimes P/\hbar$；幺正性及 $X^2$ 与传播子的对易性给出上述范数界。
初始探针满足 $\langle P\rangle_i=0$，所以归一化因子在一阶不变。
再保留二阶振幅，记 $A_w^{(2)}=\langle f|A^2|i\rangle/c_{fi}$，成功概率为
\[
 P_f=|c_{fi}|^2\left[
 1+\frac{g^2\sigma_P^2}{\hbar^2}
 (|A_w|^2-\operatorname{Re}A_w^{(2)})\right]+O(g^3).
\]
这一式由一阶振幅模平方与零阶、二阶振幅交叉项分别得到，
说明条件位移增大时还必须同步计算成功率。这里及下面的余项均以预选、后选态和探针宽度固定为前提。对任意探针算符 $\hat O$，条件平均值的一阶修正为
\[
\delta\langle O\rangle
=\frac{\ii g}{\hbar}
\left(A_w^*\langle P O\rangle_i
-A_w\langle O P\rangle_i\right)+O(g^2).
\]
对实、居中的 高斯 波包，
\[
\langle\{Q,P\}\rangle_i=0,
\qquad
\langle QP\rangle_i=\frac{\ii\hbar}{2},
\qquad
\langle PQ\rangle_i=-\frac{\ii\hbar}{2}.
\]
因此取 $O=Q$ 得
\begin{align*}
\delta\langle Q\rangle
&=\frac{\ii g}{\hbar}
\left[-\frac{\ii\hbar}{2}A_w^*
-\frac{\ii\hbar}{2}A_w\right]+O(g^2)\\
&=g\,\operatorname{Re}A_w+O(g^2).
\end{align*}
取 $O=P$ 并使用 $\langle P^2\rangle_i=\sigma_P^2$，
\begin{align*}
\delta\langle P\rangle
&=\frac{\ii g}{\hbar}(A_w^*-A_w)\sigma_P^2+O(g^2)\\
&=\frac{2g\sigma_P^2}{\hbar}\operatorname{Im}A_w+O(g^2).
\end{align*}
最小不确定 高斯 满足 $\sigma_Q\sigma_P=\hbar/2$，代入即得到\eqnrefs{eq:weak_value_gaussian_meter_shifts} 的最后一个等号。
\end{exbox}

\subsection{前向模型与信息界}

测量算符给出了从态到概率的规则。把相互作用和传播所得的态代入，就能得到概率对实验参数的依赖。设 $\boldsymbol\theta$ 是这些待估参数组成的向量，并用 $\mathscr C_{\rm int}$、$\mathscr C_{\rm prop}$、$\mathscr C_{\rm env}$ 分别表示相互作用、传播及环境阶段的完全正保迹映射。在这三个阶段可以按顺序分开描述的条件下，探测器前的态为
\begin{equation}
\boxed{
\rho_{\rm det}(\boldsymbol\theta)
=(\mathscr C_{\rm env}\circ\mathscr C_{\rm prop}\circ\mathscr C_{\rm int})(\boldsymbol\theta)[\rho_0].}
\label{eq:detector_pre_state_channel_composition}
\end{equation}
若理想量子结果 $x$ 的 POVM 效应算符是 $\Pi_x$，其后还有纯经典响应 $R(y|x)$，则完整实验似然函数是
\begin{equation}
\boxed{
p_{\rm meas}(y|\boldsymbol\theta)
=\sum_xR(y|x)
\operatorname{Tr}\!\left[
\Pi_x(\boldsymbol\theta_{\rm det})
(\mathscr C_{\rm env}\circ\mathscr C_{\rm prop}\circ\mathscr C_{\rm int})(\boldsymbol\theta)[\rho_0]
\right].}
\label{eq:general_experimental_forward_model}
\end{equation}
连续结果时把求和替换成对应测度积分。这个分层保证“相互作用改变状态”“未观测环境导致量子信道”“探测器把量子结果映射成数据”三件事不会被混在同一个唯象展宽中。

若目标是重建有限维子空间中的未知态，选择 Hilbert--Schmidt 正交的厄米基 $\{B_a\}$。未知态和每个设置下的 Born 概率先写成
\begin{equation}
\boxed{
\rho=\sum_a r_aB_a,
\qquad
p(y|s)=\operatorname{Tr}(E_{y|s}\rho).}
\label{eq:finite_dimensional_state_operator_expansion}
\end{equation}
其中全部已知控制都吸收到有效效应算符 $E_{y|s}$。代入基展开后，Born 规则成为线性算符框架
\begin{equation}
\boxed{
\mathbf p=\mathsf A\mathbf r,
\qquad
A_{(y,s),a}=\operatorname{Tr}(E_{y|s}B_a).}
\label{eq:measurement_operator_frame}
\end{equation}
在迹已固定为一的无迹厄米子空间上，信息完备性等价于 $\mathsf A$ 没有非零零方向：
\[
\ker\mathsf A\cap\mathsf{Herm}_0(\mathcal H_S)=\{0\}.
\]
这只是唯一性条件；实际稳定性由框架算符 $\mathsf F=\mathsf A^{\mathsf T}\mathsf W\mathsf A$ 的小奇异值控制。正则化可以抑制噪声，但不能创造测量从未编码的信息。

\begin{proposition}[有限维量子态的测量算符框架]
设 $B_a$ 为 厄米 算符空间的 Hilbert--Schmidt 正交基，$B_0=I/\sqrt d$。Born 概率满足\eqnrefs{eq:measurement_operator_frame}。固定迹后，测量在所有态上信息完备当且仅当设计矩阵在无迹 厄米 子空间上单射；正权框架矩阵的最小特征值控制逆问题稳定性。
\end{proposition}
\begin{proof}
设系统 Hilbert 空间维数为 $d$，选取 Hilbert--Schmidt 正交的厄米算符基
\[
\Tr(B_aB_b)=\delta_{ab},
\qquad a,b=0,\ldots,d^2-1.
\]
取
\[
B_0=\frac{\mathbb I}{\sqrt d},
\]
其余 $B_a$ 为无迹厄米算符。正文的状态展开
\[
\rho=\sum_a r_aB_a
\]
两边左乘 $B_b$ 并取迹：
\begin{align*}
\Tr(B_b\rho)
&=\sum_a r_a\Tr(B_bB_a)\\
&=\sum_a r_a\delta_{ba}\\
&=r_b.
\end{align*}
因此 $r_a=\Tr(B_a\rho)$，且迹条件固定
\[
r_0=\Tr(B_0\rho)=\frac1{\sqrt d}.
\]

对实验设置 $s$ 和结果 $y$，全部已知控制、传播与探测已经包含在系统侧效应 $E_{y|s}$ 中。Born 规则逐步给出
\begin{align*}
p(y|s)
&=\Tr(E_{y|s}\rho)\\
&=\Tr\!\left[E_{y|s}\sum_a r_aB_a\right]\\
&=\sum_a r_a\Tr(E_{y|s}B_a).
\end{align*}
把复合数据指标记为 $\mu=(y,s)$，定义矩阵元
\[
A_{\mu a}\equiv\Tr(E_\mu B_a),
\qquad
p_\mu\equiv p(y|s).
\]
逐个结果的标量等式
\[
p_\mu=\sum_aA_{\mu a}r_a
\]
合并成矩阵形式
\[
\mathbf p=\mathsf A\mathbf r,
\qquad
A_{(y,s),a}=\Tr(E_{y|s}B_a),
\]
即正文\eqnrefs{eq:measurement_operator_frame}。
\end{proof}


若目标只是估计参数 $\theta$，则由真实前向模型产生的概率 $p_x(\theta)$ 定义经典 Fisher 信息
\begin{equation}
\boxed{
F_C(\theta)=\sum_x\frac{[\partial_\theta p_x(\theta)]^2}{p_x(\theta)}.}
\label{eq:classical_fisher_information_definition}
\end{equation}
概率对参数的相对敏感度
\begin{equation*}
s_\theta(x)\equiv\partial_\theta\ln p_x(\theta)
\end{equation*}
称为该结果的\emph{得分函数}（score）。它的均值为零，而其平方的概率平均正是 $F_C=\sum_xp_xs_\theta^2$；因此 Fisher 信息衡量小幅改变参数时，整套观测概率分布平均能够改变多快。对多参数情形，相应得分向量为 $\mathbf s(x)=\nabla_{\boldsymbol\theta}\ln p_x$。

$N$ 次独立重复的局域无偏估计满足
\begin{equation}
\boxed{
\operatorname{Var}(\hat\theta)\ge\frac1{N F_C(\theta)}.}
\label{eq:classical_cramer_rao_general}
\end{equation}
探测器损失与模糊必须先进入 $p_x$ 再计算 Fisher 信息，因此它们不能在信息分析以后通过一个事后“校正因子”恢复。

\begin{theorem}[正则单参数 Cramer--Rao 界]
设参数无关支撑上的概率族可微，求导与求和可交换，$0<F_C<\infty$。对在参数点局部无偏且期望导数为一的估计量，$N$ 次独立同分布观测满足 $\operatorname{Var}(\hat\theta)\ge1/(NF_C)$。参数依赖支撑或不满足求导交换条件时不能直接套用此式。
\end{theorem}
\begin{proof}
设单次观测 $x$ 的概率为 $p_x(\theta)$，并定义得分函数
\[
s_\theta(x)=\partial_\theta\ln p_x(\theta)
=\frac{\partial_\theta p_x}{p_x}.
\]
先验证得分的均值为零。由概率归一化
\[
\sum_xp_x(\theta)=1
\]
对 $\theta$ 求导：
\begin{align*}
0
&=\partial_\theta\sum_xp_x(\theta)\\
&=\sum_x\partial_\theta p_x(\theta)\\
&=\sum_xp_x(\theta)s_\theta(x)\\
&=\mathbb E_\theta[s_\theta].
\end{align*}
于是 Fisher 信息就是得分方差：
\[
F_C(\theta)
=\mathbb E_\theta[s_\theta^2]
=\sum_x\frac{[\partial_\theta p_x]^2}{p_x}.
\]

令 $\hat\theta(x)$ 是局域无偏估计量，满足
\[
\mathbb E_\theta[\hat\theta]=\theta.
\]
对该等式求导：
\begin{align*}
1
&=\partial_\theta\sum_xp_x(\theta)\hat\theta(x)\\
&=\sum_x\hat\theta(x)\partial_\theta p_x(\theta)\\
&=\sum_xp_x(\theta)\hat\theta(x)s_\theta(x).
\end{align*}
因为 $\mathbb E[s_\theta]=0$，可减去任意常数，特别取真参数 $\theta$：
\begin{align*}
1
&=\sum_xp_x(\theta)[\hat\theta(x)-\theta]s_\theta(x)\\
&=\mathbb E_\theta[(\hat\theta-\theta)s_\theta].
\end{align*}
对随机变量 $\hat\theta-\theta$ 与 $s_\theta$ 使用 Cauchy--Schwarz 不等式：
\begin{align*}
1
&\le\sqrt{\mathbb E[(\hat\theta-\theta)^2]}
\sqrt{\mathbb E[s_\theta^2]}\\
&=\sqrt{\operatorname{Var}(\hat\theta)F_C(\theta)}.
\end{align*}
平方并整理得到单次观测的界
\[
\operatorname{Var}(\hat\theta)\ge\frac1{F_C(\theta)}.
\]

对 $N$ 次相互独立、同分布观测
$\mathbf x=(x_1,\ldots,x_N)$，联合概率为
\[
p(\mathbf x|\theta)=\prod_{j=1}^Np_{x_j}(\theta),
\]
所以联合得分是各次得分之和：
\[
s_\theta^{(N)}(\mathbf x)
=\partial_\theta\ln p(\mathbf x|\theta)
=\sum_{j=1}^Ns_\theta(x_j).
\]
独立性和零均值使交叉项消失：
\begin{align*}
F_C^{(N)}
&=\mathbb E[(s_\theta^{(N)})^2]\\
&=\sum_{j=1}^N\mathbb E[s_\theta^2(x_j)]
+\sum_{j\ne k}\mathbb E[s_\theta(x_j)]\mathbb E[s_\theta(x_k)]\\
&=NF_C(\theta).
\end{align*}
代回单次形式即得
\[
\operatorname{Var}(\hat\theta)
\ge\frac1{NF_C(\theta)},
\]
即正文\eqnrefs{eq:classical_cramer_rao_general}。
\end{proof}


真实反演通常同时包含多个未知参数。对参数向量 $\boldsymbol\theta=(\theta_1,\ldots,\theta_d)$，经典 Fisher 信息推广为矩阵
\begin{equation}
 \boxed{
 [\mathsf I_C(\boldsymbol\theta)]_{ab}
 =\sum_x\frac{\partial_a p_x\,\partial_b p_x}{p_x}.}
 \label{eq:multiparameter-fisher-matrix-dp31}
\end{equation}
若参数分成目标参数 $\psi$ 与干扰参数向量 $\boldsymbol\theta_{\rm n}$，并且干扰参数块 $\mathsf I_{{\rm n}{\rm n}}$ 可逆，则在同时拟合 $\boldsymbol\theta_{\rm n}$ 时，$\psi$ 的局部有效信息是 Schur 补
\begin{equation}
 \boxed{
 I_{\rm eff}(\psi)
 =I_{\psi\psi}
 -\mathsf I_{\psi{\rm n}}
 \mathsf I_{{\rm n}{\rm n}}^{-1}
 \mathsf I_{{\rm n}\psi}.}
 \label{eq:fisher-schur-nuisance-dp31}
\end{equation}
第二项非负，因此不知道干扰参数只会降低目标参数的可辨识性。若该 Schur 补为零，即使 $I_{\psi\psi}$ 本身很大，目标参数也会因为与干扰方向局部退化而不可辨识。

\begin{proposition}[干扰得分投影与有效 Fisher 信息]
设正则多参数 Fisher 矩阵正定，目标参数为标量，干扰块 $\mathsf C$ 可逆。则有效信息为 $a-\mathbf b^{\mathsf T}\mathsf C^{-1}\mathbf b$，等于目标得分 减去干扰得分 正交投影后的平方范数；局部无偏估计的目标方差不低于其倒数除以 $N$，即\eqnrefs{eq:fisher-schur-nuisance-dp31}。
\end{proposition}
\begin{proof}
把待估参数拆成目标标量 $\psi$ 与干扰参数向量 $\boldsymbol\theta_{\rm n}$。单次实验的 Fisher 矩阵写成
\[
\mathsf I_C=
\begin{pmatrix}
 a & \mathbf b^{\mathsf T}\\
 \mathbf b & \mathsf C
\end{pmatrix},
\qquad
 a\equiv I_{\psi\psi},\quad
 \mathbf b\equiv\mathsf I_{{\rm n}\psi},\quad
 \mathsf C\equiv\mathsf I_{{\rm n}{\rm n}}.
\]
先假设 $\mathsf C$ 在局部可辨识的干扰子空间上可逆。与其直接引用分块逆公式，不妨解线性方程
\[
\mathsf I_C
\begin{pmatrix}x\\ \mathbf y\end{pmatrix}
=
\begin{pmatrix}1\\ \mathbf 0\end{pmatrix}.
\]
第二行给
\[
\mathbf b\,x+\mathsf C\mathbf y=0
\quad\Longrightarrow\quad
\mathbf y=-\mathsf C^{-1}\mathbf b\,x.
\]
代回第一行：
\begin{align*}
ax+\mathbf b^{\mathsf T}\mathbf y&=1,\\
\left(a-\mathbf b^{\mathsf T}\mathsf C^{-1}\mathbf b\right)x&=1.
\end{align*}
而 $x$ 正是 $\mathsf I_C^{-1}$ 的左上角元素，因此
\[
[\mathsf I_C^{-1}]_{\psi\psi}
=\frac{1}{a-\mathbf b^{\mathsf T}\mathsf C^{-1}\mathbf b}.
\]
定义
\[
\boxed{
I_{\rm eff}(\psi)
\equiv
I_{\psi\psi}
-\mathsf I_{\psi{\rm n}}
\mathsf I_{{\rm n}{\rm n}}^{-1}
\mathsf I_{{\rm n}\psi}},
\]
多参数 Cram\'er--Rao 界
\[
\operatorname{Cov}(\hat{\boldsymbol\theta})
\succeq\frac1N\mathsf I_C^{-1}
\]
就在目标参数方向给出
\[
\boxed{
\operatorname{Var}(\hat\psi)
\ge\frac1{N I_{\rm eff}(\psi)}}.
\]

扣除项的来源也可从 score 几何直接看出。记单次数据的目标得分 为 $s_\psi$，干扰得分 向量为 $\mathbf s_{\rm n}$，则
\[
a=\mathbb E[s_\psi^2],\qquad
\mathbf b=\mathbb E[\mathbf s_{\rm n}s_\psi],\qquad
\mathsf C=\mathbb E[\mathbf s_{\rm n}\mathbf s_{\rm n}^{\mathsf T}].
\]
在线性组合 $s_\psi-\mathbf c^{\mathsf T}\mathbf s_{\rm n}$ 中选择 $\mathbf c$ 使均方最小：
\[
\frac{\partial}{\partial\mathbf c}
\mathbb E[(s_\psi-\mathbf c^{\mathsf T}\mathbf s_{\rm n})^2]
=-2\mathbf b+2\mathsf C\mathbf c=0,
\]
故 $\mathbf c_*=\mathsf C^{-1}\mathbf b$，最小剩余范数恰为
\[
\mathbb E[(s_\psi-\mathbf c_*^{\mathsf T}\mathbf s_{\rm n})^2]
=a-\mathbf b^{\mathsf T}\mathsf C^{-1}\mathbf b
=I_{\rm eff}.
\]
所以 Schur 补不是抽象的矩阵技巧：它把目标得分 中能够由干扰参数 score 线性模仿的方向投影掉，只保留实验真正能独立归因于 $\psi$ 的信息。若 $\mathsf C$ 奇异，应限制到其可辨识支撑空间或使用 Moore--Penrose 逆；零模对应数据无法区分的真实参数退化。
\end{proof}


纯经典、且与待估参数无关的探测器响应也满足信息数据处理不等式。若理想结果 $x$ 经过随机核 $R(y|x)$ 变成记录 $y$，粗粒化概率满足
\begin{equation}
\boxed{
\bar p_y(\boldsymbol\theta)=\sum_xR(y|x)p_x(\boldsymbol\theta),
\qquad
\sum_yR(y|x)=1.}
\label{eq:classical_detector_coarse_graining}
\end{equation}
对这一概率重新计算 Fisher 矩阵，得到
\begin{equation}
 \boxed{
 \mathsf I_C^{(y)}\preceq\mathsf I_C^{(x)}.}
 \label{eq:fisher-data-processing-dp31}
\end{equation}
矩阵差为半正定。能量展宽、像素合并、漏计数后的条件重标记等只要可表示成参数无关的经典粗粒化，就不能创造新的 Fisher 信息；若探测器本身含未知参数，则这些参数必须并入 $\boldsymbol\theta$，再用上式的多参数结构共同估计。

\begin{theorem}[参数无关后处理的 Fisher 收缩]
设正则概率族经参数无关 Markov 核 $R(y|x)$ 后处理。则记录 score 等于理想 score 的条件期望，因而 $\mathsf I_C^{(y)}\preceq\mathsf I_C^{(x)}$。等号当且仅当全部 score 方向几乎处处可由记录 $Y$ 确定。
\end{theorem}
\begin{proof}
设理想量子结果为 $x$，实际记录为 $y$。纯经典探测器后处理由与待估参数 $\boldsymbol\theta$ 无关的 Markov 核 $R(y|x)$ 描述：
\[
R(y|x)\ge0,
\qquad
\sum_yR(y|x)=1,
\qquad
\bar p_y=\sum_xR(y|x)p_x.
\]
对参数求导时，$R$ 本身没有导数，因此
\begin{align*}
\nabla_{\boldsymbol\theta}\bar p_y
&=\sum_xR(y|x)\nabla_{\boldsymbol\theta}p_x\\
&=\sum_xR(y|x)p_x\,\mathbf s(x),
\end{align*}
其中理想数据的 score 向量定义为
\[
\mathbf s(x)=\nabla_{\boldsymbol\theta}\ln p_x.
\]
实际记录 $y$ 的 score 为
\begin{align*}
\bar{\mathbf s}(y)
&=\nabla_{\boldsymbol\theta}\ln\bar p_y\\
&=\frac1{\bar p_y}
\sum_xR(y|x)p_x\mathbf s(x).
\end{align*}
定义给定 $y$ 后理想结果 $x$ 的后验权重
\[
w(x|y)=\frac{R(y|x)p_x}{\bar p_y},
\qquad
\sum_xw(x|y)=1.
\]
于是
\[
\bar{\mathbf s}(y)=\sum_xw(x|y)\mathbf s(x)
=\mathbb E[\mathbf s(X)|Y=y].
\]
为了证明 Fisher 矩阵的半正定序关系，取任意实参数方向 $\mathbf v$，并定义标量投影
\[
s_x=\mathbf v^{\mathsf T}\mathbf s(x),
\qquad
t_y=\mathbf v^{\mathsf T}\bar{\mathbf s}(y).
\]
由上式
\[
t_y=\mathbb E[s_X|Y=y].
\]
对凸函数 $f(s)=s^2$ 使用条件 Jensen 不等式：
\[
t_y^2
=\bigl(\mathbb E[s_X|Y=y]\bigr)^2
\le\mathbb E[s_X^2|Y=y].
\]
乘以 $\bar p_y$ 并对 $y$ 求和：
\begin{align*}
\sum_y\bar p_yt_y^2
&\le\sum_y\bar p_y
\sum_xw(x|y)s_x^2\\
&=\sum_{x,y}R(y|x)p_xs_x^2\\
&=\sum_xp_xs_x^2\sum_yR(y|x)\\
&=\sum_xp_xs_x^2.
\end{align*}
两端分别就是实际记录和理想结果的 Fisher 二次型：
\begin{align*}
\mathbf v^{\mathsf T}\mathsf I_C^{(y)}\mathbf v
&=\sum_y\bar p_y
[\mathbf v^{\mathsf T}\bar{\mathbf s}(y)]^2,\\
\mathbf v^{\mathsf T}\mathsf I_C^{(x)}\mathbf v
&=\sum_xp_x
[\mathbf v^{\mathsf T}\mathbf s(x)]^2.
\end{align*}
因此对任意 $\mathbf v$ 都有
\[
\mathbf v^{\mathsf T}
(\mathsf I_C^{(x)}-\mathsf I_C^{(y)})\mathbf v\ge0,
\]
这正是矩阵半正定序
\[
\mathsf I_C^{(y)}\preceq\mathsf I_C^{(x)},
\]
即正文\eqnrefs{eq:fisher-data-processing-dp31}。
\end{proof}


量子 Fisher 信息则是给定态族 $\rho_\theta$ 时对所有 POVM 的共同上界。对称对数导数 $L_\theta$ 与量子 Fisher 信息定义为
\begin{equation}
\boxed{
\partial_\theta\rho_\theta
=\frac12(L_\theta\rho_\theta+\rho_\theta L_\theta),
\qquad
F_Q=\operatorname{Tr}(\rho_\theta L_\theta^2).}
\label{eq:qfi-sld-definition-dp68}
\end{equation}
对任意 POVM，经典 Fisher 信息满足
\begin{equation}
\boxed{F_C\le F_Q.}
\label{eq:qfi-measurement-upper-bound-dp68}
\end{equation}
这个不等式把“态本身携带的局域参数信息”与“某个具体测量实际读出的信息”区分开来。对纯态，
\begin{equation}
\boxed{
F_Q
=4\left[
\langle\partial_\theta\Psi|\partial_\theta\Psi\rangle
-|\langle\Psi|\partial_\theta\Psi\rangle|^2
\right].}
\label{eq:pure_state_quantum_fisher_general}
\end{equation}
\eqnrefs{eq:pure_state_quantum_fisher_general} 中第二项恰好扣除了 $|\Psi_\theta\rangle$ 方向上的纯相位变化，因此纯态 QFI 只保留真正改变量子态射线的正交分量。若参数由无量纲厄米生成元 $G$ 编码，$|\Psi_\theta\rangle=\ee^{-\ii\theta G}|\Psi_0\rangle$，则
\begin{equation}
\boxed{F_Q=4(\Delta G)^2.}
\label{eq:unitary_encoding_qfi_variance}
\end{equation}
若实际参数为带量纲的 $\lambda$ 且编码算符为 $\exp[-\ii\lambda K/\hbar]$，则 $G=K/\hbar$，从而 $F_Q^{(\lambda)}=4(\Delta K)^2/\hbar^2$；若生成元本身依赖参数，则必须对指数算符使用有序积分形式的导数，不能直接套用方差公式。因此增强耦合并不自动等价于更高可辨识性；真正的问题是态族怎样随参数改变、实际 POVM 读出了其中多少信息，以及噪声信道是否把这些变化投影到了不可观测方向。

\begin{proposition}[纯态 QFI 与幺正生成元方差]
设归一化纯态曲线在 Hilbert 范数中可微，则 $F_Q=4(\langle\dot\Psi|\dot\Psi\rangle-|\langle\Psi|\dot\Psi\rangle|^2)$。若 $|\Psi_\theta\rangle=e^{-\ii\theta G}|\Psi_0\rangle$，$G$ 自伴且初态有有限二阶矩，则 $F_Q=4\operatorname{Var}(G)$。
\end{proposition}
\begin{proof}
令
\begin{equation*}
\rho=|\Psi\rangle\langle\Psi|,
\qquad
|\dot\Psi\rangle\equiv\partial_\theta|\Psi_\theta\rangle.
\end{equation*}
归一化条件给
\begin{equation*}
0=\partial_\theta\langle\Psi|\Psi\rangle
=\langle\dot\Psi|\Psi\rangle+\langle\Psi|\dot\Psi\rangle.
\end{equation*}
记 $a=\langle\Psi|\dot\Psi\rangle$，并分解
\begin{equation*}
|\dot\Psi\rangle=a|\Psi\rangle+|\dot\Psi_\perp\rangle,
\qquad
\langle\Psi|\dot\Psi_\perp\rangle=0.
\end{equation*}
由于 $a+a^*=0$，
\begin{equation*}
\dot\rho
=|\dot\Psi_\perp\rangle\langle\Psi|
+|\Psi\rangle\langle\dot\Psi_\perp|.
\end{equation*}
取
\begin{equation*}
L=2\left(
|\dot\Psi_\perp\rangle\langle\Psi|
+|\Psi\rangle\langle\dot\Psi_\perp|
\right),
\end{equation*}
则利用 $\rho|\dot\Psi_\perp\rangle=0$ 可验证
\begin{equation*}
\frac12(L\rho+\rho L)=\dot\rho.
\end{equation*}
再令 $|v\rangle=|\dot\Psi_\perp\rangle$，有
\begin{equation*}
\frac{L^2}{4}
=|v\rangle\langle v|+\langle v|v\rangle|\Psi\rangle\langle\Psi|,
\end{equation*}
因此
\begin{align*}
F_Q
&=\operatorname{Tr}(\rho L^2)
=4\langle v|v\rangle\\
&=4\left[
\langle\dot\Psi|\dot\Psi\rangle
-|\langle\Psi|\dot\Psi\rangle|^2
\right],
\end{align*}
即正文\eqnrefs{eq:pure_state_quantum_fisher_general}。

若 $|\Psi_\theta\rangle=\ee^{-\ii\theta G}|\Psi_0\rangle$ 且 $G=G^\dagger$ 与 $\theta$ 无关，则
\begin{equation*}
|\dot\Psi_\theta\rangle=-\ii G|\Psi_\theta\rangle,
\end{equation*}
从而
\begin{equation*}
\langle\dot\Psi|\dot\Psi\rangle=\langle G^2\rangle,
\qquad
|\langle\Psi|\dot\Psi\rangle|^2=\langle G\rangle^2.
\end{equation*}
代入上一式得到
\begin{equation*}
F_Q=4\left(\langle G^2\rangle-\langle G\rangle^2\right)
=4(\Delta G)^2,
\end{equation*}
即正文\eqnrefs{eq:unitary_encoding_qfi_variance}。
\end{proof}


\begin{theorem}[SLD 量子信息对测量信息的上界]
设有限维可微态族存在 厄米 SLD $L$，$\partial_\theta\rho=(L\rho+\rho L)/2$。对参数无关 POVM，在参数点的正概率支撑上定义经典 Fisher 信息，则 $F_C\le\Tr(\rho L^2)=F_Q$。常秩正则模型中零概率项按支撑处理；跨秩边界的极限 Fisher 信息须单独分析。
\end{theorem}
\begin{proof}
设 POVM 元为 $E_x\ge0$、$\sum_xE_x=I$，测量概率为
$p_x=\operatorname{Tr}(E_x\rho_\theta)$。SLD 定义给
\begin{align*}
 \partial_\theta p_x
 &=\operatorname{Tr}(E_x\partial_\theta\rho_\theta)\\
 &=\frac12\operatorname{Tr}[E_x(L\rho+\rho L)]\\
 &=\operatorname{Re}\operatorname{Tr}(E_xL\rho).
\end{align*}
最后一步使用 $E_x,L,\rho$ 的厄米性：两项互为复共轭。

为了把每个概率导数与 QFI 中的 $L^2$ 联系起来，在 Hilbert--Schmidt 内积
$\langle A,B\rangle_{\rm HS}=\operatorname{Tr}(A^\dagger B)$ 下定义
\begin{equation*}
 A_x=E_x^{1/2}\rho^{1/2},
 \qquad
 B_x=E_x^{1/2}L\rho^{1/2}.
\end{equation*}
于是
\begin{align*}
 \operatorname{Tr}(A_x^\dagger B_x)
 &=\operatorname{Tr}(\rho^{1/2}E_xL\rho^{1/2})
 =\operatorname{Tr}(E_xL\rho),\\
 \operatorname{Tr}(A_x^\dagger A_x)
 &=\operatorname{Tr}(E_x\rho)=p_x,\\
 \operatorname{Tr}(B_x^\dagger B_x)
 &=\operatorname{Tr}(\rho L E_xL).
\end{align*}
因此先用 $|\operatorname{Re}z|^2\le|z|^2$，再用 Cauchy--Schwarz，得到
\begin{equation*}
 |\partial_\theta p_x|^2
 \le |\operatorname{Tr}(A_x^\dagger B_x)|^2
 \le p_x\operatorname{Tr}(\rho L E_xL).
\end{equation*}
先处理 $p_x=0$ 的边界情形。因为 $E_x\ge0$ 与 $\rho\ge0$，
\begin{equation*}
p_x=\operatorname{Tr}(E_x\rho)
=\operatorname{Tr}\!\left[
(E_x^{1/2}\rho^{1/2})^\dagger
(E_x^{1/2}\rho^{1/2})
\right]=0
\end{equation*}
蕴含
\begin{equation*}
E_x^{1/2}\rho^{1/2}=A_x=0.
\end{equation*}
于是
$\partial_\theta p_x=\operatorname{Re}\operatorname{Tr}(A_x^\dagger B_x)=0$；在参数点附近可微的正则模型中，这个结果对 Fisher 和式不产生奇异贡献。因而只需对 $p_x>0$ 的支撑除以 $p_x$：
\begin{equation*}
\frac{(\partial_\theta p_x)^2}{p_x}
\le\operatorname{Tr}(\rho L E_xL).
\end{equation*}
随后对所有测量结果求和，并使用 POVM 完备性：
\begin{align*}
 F_C
 &=\sum_{x:p_x>0}\frac{(\partial_\theta p_x)^2}{p_x}\\
 &\le\sum_x\operatorname{Tr}(\rho L E_xL)\\
 &=\operatorname{Tr}\!\left[
\rho L\left(\sum_xE_x\right)L
\right]\\
 &=\operatorname{Tr}(\rho L^2)\\
 &=F_Q.
\end{align*}
从概率导数到最后一行只用了 SLD 定义、Hilbert--Schmidt Cauchy--Schwarz 不等式和 $\sum_xE_x=I$，因此该界与具体探测器实现无关。这就是正文\eqnrefs{eq:qfi-measurement-upper-bound-dp68}。
\end{proof}


等号要求 Cauchy--Schwarz 与取实部两步同时饱和；局域测量 SLD 本征投影可满足这一结构。因此 QFI 是对允许 POVM 优化后的局域上界，而不是固定探测器自动达到的灵敏度。

如果实验目标不是重建单一时刻的电子态，而是识别多电子连续碰撞中的环境记忆，那么未知对象应从密度矩阵升级为前文定义的过程张量 $\Upsilon_{N:0}$。这时“数据足够多”必须有一个精确的线性代数含义。因果归一化含有非零恒等项，因此允许过程构成仿射空间。选取一个固定参考过程 $\Upsilon_0$，再取齐次因果约束下的独立厄米方向 $\{F_\alpha\}_{\alpha=1}^{d_{\rm proc}}$，写
\begin{equation*}
\Upsilon_{N:0}=\Upsilon_0+\sum_{\alpha=1}^{d_{\rm proc}}\theta_\alpha F_\alpha.
\end{equation*}
一次实验设置 $r$ 由各时刻插入的电子制备/控制 Choi 算符
$A_{r,0},\ldots,A_{r,N-1}$ 和末端 POVM 元 $M_{r,y}$ 决定。连接积 对每个槽位都是线性的，因此结果概率必可写成
\begin{equation}
\boxed{
p_{r,y}
=p^{(0)}_{r,y}+\sum_{\alpha=1}^{d_{\rm proc}}
X_{(r,y),\alpha}\theta_\alpha,
\qquad
X_{(r,y),\alpha}
\equiv
\Tr\!\left[
M_{r,y}
\left(F_\alpha\star A_{r,N-1}\star\cdots\star A_{r,0}\right)
\right].}
\label{eq:process_tensor_experimental_design_matrix}
\end{equation}
这里 $p^{(0)}_{r,y}$ 是对参考过程实施同一实验得到的概率，$X$ 只描述齐次允许方向。对每个完整结果集合，$\sum_y p^{(0)}_{r,y}=1$ 且 $\sum_yX_{(r,y),\alpha}=0$；过程的正性还限制可行参数。过程层析能否唯一完成，不取决于重复次数本身，而取决于实验设计矩阵 $X$ 是否覆盖所有允许的过程方向：
\begin{equation}
\boxed{
\operatorname{rank}X=d_{\rm proc}.}
\label{eq:process_tensor_identifiability_rank}
\end{equation}
若秩亏，增加同一设置的重复计数只能降低已可见方向的统计噪声，不能恢复落在 $\ker X$ 中的记忆参数。

对自由电子实验，这个结论具有直接操作意义。若每个电子都以完全相同的波包入射、使用同一相互作用强度、相同电子间隔，并且最终只测一个固定能谱，那么实验只探测过程张量的一个很低维切片。要区分“电子自身相干”“环境跨碰撞记忆”“材料慢漂移”和“探测器卷积”，通常必须改变至少一类控制：入射电子相位/能量基、电子间延迟、耦合强度、后选择基或末端探测基。过程张量形式把这些变化视为不同的槽位操作，而不是不同的理论模型。

在大样本局部估计中，若计数服从独立多项式/Poisson 统计，过程参数的经典 Fisher 信息仍由概率导数组成：
\begin{equation}
\boxed{
[\mathcal F_{\rm proc}]_{\alpha\beta}
=\sum_{r,y}
\frac{N_r}{p_{r,y}}
X_{(r,y),\alpha}X_{(r,y),\beta},}
\label{eq:process_tensor_fisher_information}
\end{equation}
其中 $N_r$ 是设置 $r$ 的重复次数。若所有被使用的 $p_{r,y}>0$，则
$\mathcal F_{\rm proc}$ 在允许参数空间上正定与 $X$ 满列秩等价。因此“过程可辨识性”和“有限方差估计”由同一个设计矩阵控制；实验优化可以转化为改善 $X$ 的奇异值谱，而不是只追求更多电子计数。

\begin{proposition}[过程设计矩阵的 Fisher 零空间]
设过程参数采用允许仿射空间的独立实坐标，概率满足 $p=p^{(0)}+X\theta$，使用的 $p_{r,y}>0$ 且重复数 $N_r>0$。独立多项式计数，或已知曝光 $N_r$ 下均值 $N_rp_{r,y}$ 的独立 Poisson 计数，都给\eqnrefs{eq:process_tensor_fisher_information}，并满足 $\ker\mathcal F_{\rm proc}=\ker X$。物理集合边界上的可行方向须另行限制。
\end{proposition}
\begin{proof}
参考过程承担非齐次归一化，变化方向 $F_\alpha$ 承担齐次因果约束，故过程 Choi 算符展开为
\[
\Upsilon=\Upsilon_0+\sum_\alpha\theta_\alpha F_\alpha.
\]
固定一个实验设置 $r$。多时刻操作的 连接积 定义一个从过程算符到末态的线性映射
\[
\mathcal L_r(F)
\equiv
F\star A_{r,N-1}\star\cdots\star A_{r,0}.
\]
末端 Born 规则给
\begin{align*}
p_{r,y}
&=\Tr[M_{r,y}\mathcal L_r(\Upsilon)]\\
&=\Tr[M_{r,y}\mathcal L_r(\Upsilon_0)]
+\sum_\alpha\theta_\alpha\Tr[M_{r,y}\mathcal L_r(F_\alpha)]\\
&\equiv p^{(0)}_{r,y}+\sum_\alpha X_{(r,y),\alpha}\theta_\alpha,
\end{align*}
即正文\eqnrefs{eq:process_tensor_experimental_design_matrix}。

若存在非零向量 $v\in\ker X$，则对任意足够小的实数 $\epsilon$，参数
$\theta$ 与 $\theta+\epsilon v$ 产生完全相同的全部一阶概率变化；若二者都保持在物理过程集合内，数据无法区分它们。因此 $\ker X\neq\{0\}$ 意味着局部不可辨识。反之，若 $X$ 满列秩，则线性方程组在无噪声极限中至多有一个参数解，所以局部线性层析唯一；得到正文\eqnrefs{eq:process_tensor_identifiability_rank}。

对独立计数，设置 $r$ 重复 $N_r$ 次。参数导数为
\[
\partial_{\theta_\alpha}p_{r,y}=X_{(r,y),\alpha}.
\]
经典 Fisher 信息定义
\[
[\mathcal F]_{\alpha\beta}
=\sum_{r,y}N_r p_{r,y}
(\partial_\alpha\ln p_{r,y})
(\partial_\beta\ln p_{r,y}),
\]
于是
\[
[\mathcal F]_{\alpha\beta}
=\sum_{r,y}
\frac{N_r}{p_{r,y}}
X_{(r,y),\alpha}X_{(r,y),\beta},
\]
得到正文\eqnrefs{eq:process_tensor_fisher_information}。对任意参数方向 $v$，
\[
v^{\mathsf T}\mathcal Fv
=\sum_{r,y}\frac{N_r}{p_{r,y}}
|(Xv)_{r,y}|^2.
\]
在 $p_{r,y}>0$ 的采样支撑上，这个量为零当且仅当 $Xv=0$。所以 Fisher 信息的零模正是设计矩阵不可见的过程方向。增大 $N_r$ 只放大已有非零奇异方向的权重，不能改变 $\ker X$；这严格证明了为什么“无限重复同一种电子实验”仍不能完成一般非 Markov 过程层析。
\end{proof}


% DeepPass22: detailed pure-state QFI and SLD algebra.



\begin{definition}[Poisson 与固定总数的计数模型]
设入射总数服从均值为 $N$ 的 Poisson 分布，$N$ 已知；各电子独立进入互斥的探测通道，
前向模型给通道概率 $p_x(\boldsymbol\vartheta)$。Poisson 分流使通道计数相互独立，
均值为 $\lambda_x=Np_x$。其联合概率与对数似然为
\begin{equation}
 P(\mathbf n|\boldsymbol\vartheta)=\prod_x\frac{\ee^{-\lambda_x}\lambda_x^{n_x}}{n_x!},\qquad
 \log L=\sum_x[n_x\log\lambda_x-\lambda_x-\log(n_x!)].
 \label{eq:feqo-count-likelihood}
\end{equation}
得分和信息矩阵是
\begin{equation}
 S_a=\partial_a\log L=\sum_x(n_x/\lambda_x-1)\partial_a\lambda_x,
 \qquad\mathcal F_{ab}=\sum_x\frac{\partial_a\lambda_x\partial_b\lambda_x}{\lambda_x}.
 \label{eq:feqo-poisson-fisher}
\end{equation}
这由 $\mathbb E n_x=\lambda_x$、$\operatorname{Cov}(n_x,n_y)=\delta_{xy}\lambda_x$ 得到。
若固定总计数为整数 $N_0$ 且完整结果空间 $\sum_xp_x=1$，则
\begin{equation}
 P(\mathbf n|\boldsymbol\vartheta,N_0)=\frac{N_0!}{\prod_xn_x!}\prod_xp_x^{n_x},\qquad
 \sum_xn_x=N_0,\qquad
 \mathcal F_{ab}=N_0\sum_x\frac{\partial_ap_x\partial_bp_x}{p_x}.
 \label{eq:feqo-multinomial}
\end{equation}
固定入射数、固定探测成功数和仅观察一个不完整窗口，是不同的实验条件。
计数背景、未知通量和效率应作为前向模型中的参数；若同时估计它们，
有效信息应消去相应的 干扰参数，而不能继续把 $N$ 当作已知。
\end{definition}

\begin{proposition}[计数实验的 Cram\'er--Rao 界]
\label{prop:cramer-rao}
设结果概率 $P_\vartheta(x)>0$ 在参数开区间内可微，
求和与求导可交换，且得分 $S=\partial_\vartheta\log P_\vartheta$ 有有限二阶矩。
若 $T(x)$ 在该区间无偏、方差有限，信息 $I=\mathbb E S^2>0$，则
\begin{equation}
\operatorname{Var}_\vartheta T\ge I^{-1}.
\label{eq:feqo-cr-bound}
\end{equation}
\end{proposition}
\begin{proof}
由归一化求导得 $\mathbb ES=\sum_x\partial_\vartheta P_\vartheta(x)=0$。
再对 $\mathbb ET=\vartheta$ 求导：
\[
1=\sum_x T(x)\partial_\vartheta P_\vartheta(x)
=\mathbb E[(T-\vartheta)S].
\]
Cauchy--Schwarz 给 $1\le\operatorname{Var}(T)\,I$。
对多参数无偏向量，$\mathbb E[(T-\boldsymbol\vartheta)S^{\mathsf T}]=I_d$。
信息矩阵 $\mathcal F$ 可逆时，
$Z=T-\boldsymbol\vartheta-\mathcal F^{-1}S$ 的协方差为
\begin{equation}
\operatorname{Cov}(T)-\mathcal F^{-1}=\operatorname{Cov}(Z)\succeq0.
\label{eq:feqo-cr-matrix}
\end{equation}
边界参数、零概率结果和偏差估计不自动满足这些假设。
\end{proof}

\begin{exbox}[能谱估计耦合：完整窗口、漏计与后选择]
对理想低反冲 PINEM 能谱 $p_\ell(b)=J_\ell(2b)^2$，取已知平均入射计数 $N$。
求完整窗口与有限窗口的耦合信息，并比较已知入射数与固定成功计数的差别。

\textbf{解。}
独立 Poisson 通道给信息
\[
I_b=N\sum_\ell\frac{(\partial_bp_\ell)^2}{p_\ell}
=4N\sum_\ell[J_{\ell-1}(2b)-J_{\ell+1}(2b)]^2.
\]
用 Fourier 平移正交性 $\sum_\ell J_\ell(x)J_{\ell+q}(x)=\delta_{q0}$，得到
\begin{equation}
I_b=8N,\qquad \sigma_{\widehat b}\ge(8N)^{-1/2}.
\label{eq:feqo-fisher-g-result}
\end{equation}
此等式适用于所有相关通道概率非零的正则参数点；孤立 Bessel 零点可取邻近极限，
但在零点本身必须单独检查统计模型的可微性。$b=0$ 又是模长参数的边界，
不能直接套用开区间的无偏 Cram\'er--Rao 结论。
弱耦合时 $p_0=1-2b^2+O(b^4)$、$p_{\pm1}=b^2+O(b^4)$，三者的和才在此阶等于一。

若只记录集合 $\mathcal L$ 中的 Poisson 通道且入射 $N$ 已知，
\begin{equation}
I_b^{\rm win}=N\sum_{\ell\in\mathcal L}\frac{(p_\ell')^2}{p_\ell}.
\label{eq:feqo-fisher-truncated}
\end{equation}
如果还能记录窗口外总数，就应另加
$N(q')^2/(1-q)$，其中 $q=\sum_{\ell\in\mathcal L}p_\ell$。
若只条件于探测成功后的固定计数 $N_{\rm det}$，则结果概率变成
$\widetilde p_\ell=p_\ell/q$，信息变为
\[
I_b^{\rm cond}=N_{\rm det}\left[
\frac1q\sum_{\ell\in\mathcal L}\frac{(p_\ell')^2}{p_\ell}
-\left(\frac{q'}q\right)^2\right].
\]
把这些计数协议混用会把漏计率的信息误算成边带形状信息。
例如 $b=3$、$\mathcal L=\{-2,-1,0,1,2\}$ 时，
窗口 Poisson 信息是 $3.6896N$，而完整窗口是 $8N$。
\end{exbox}

\begin{exbox}[相位参考如何把相干转成 Fisher 信息]
取两通道态
\begin{equation}
\rho=\frac12\begin{pmatrix}1&r\ee^{-\ii\phi}\\r\ee^{\ii\phi}&1\end{pmatrix},
\qquad0<r<1.
\label{eq:feqo-two-level-state}
\end{equation}
求两种参考设置的可辨识性，再计算退相干后相位信息。

\textbf{解。}
两个能量概率恒为 $1/2$，所以能谱对 $r,\phi$ 均无信息。
投影到 $|+_\theta\rangle=(|0\rangle+\ee^{\ii\theta}|1\rangle)/\sqrt2$ 得
\begin{equation}
p_+(\theta)=\frac{1+r\cos(\phi-\theta)}2,
\qquad p_-=1-p_+.
\label{eq:feqo-reference-phase-probability}
\end{equation}
$\theta=0,\pi/2$ 的两个概率关于 $(r,\phi)$ 的 Jacobian 为
\begin{equation}
J=\frac12\begin{pmatrix}\cos\phi&-r\sin\phi\\
\sin\phi&r\cos\phi\end{pmatrix},\qquad\det J=r/4.
\label{eq:feqo-jacobian-two-level}
\end{equation}
因此两设置可以恢复两个局部参数。$r=0$ 时相位没有物理定义，任何参考都无法恢复它。

纯退相干使 $r\to V=rD$，但不改变能谱人口。
固定参考相位，$N$ 个二项计数给相位信息
\begin{equation}
I_{\phi\phi}(\theta)
=\frac{NV^2\sin^2(\phi-\theta)}{1-V^2\cos^2(\phi-\theta)}.
\label{eq:feqo-fisher-phase-decoherence}
\end{equation}
这个公式由 $N(\partial_\phi p_+)^2/[p_+(1-p_+)]$ 直接得到。
最敏感的正交相位点 $\phi-\theta=\pi/2$ 给 $I_{\phi\phi}=Nr^2D^2$；
同相点在 $V<1$ 时导数为零。
若 $D=\exp(-\Gamma t)$、$r$ 已知且信号本身不随时间增长，
最敏感点的方差下界是 $\exp(2\Gamma t)/(Nr^2)$，越早测量越有利。
若同时估计 $r$，单个设置只有一个独立概率，必须增加参考设置并使用完整信息矩阵。
\end{exbox}


\subsection{可扫描参考操作与层析可辨识性}

在能量基底中，能谱仪给出的概率只含密度矩阵的对角元。若还要确定非对角相干元，可以在探测前施加一个已知操作，让不同通道先发生干涉，再测各条边带的概率。改变参考操作，便得到关于同一未知密度矩阵的多组线性方程。从这些概率恢复量子态的过程称为\emph{量子态层析}。本节用线性代数判断这些方程是否足以确定所需态分量。

\begin{exbox}[能量格与纵向相空间层析]
电子能谱仪只直接读取电子密度算符在能量基底中的对角元；若未知态含有不同能量格点之间的相干，单次能量测量不足以确定这些非对角元。层析因此需要一个\emph{已知、可重复且可扫描}的参考操作，把不同初始能量路径相干地汇合到同一个末态能量通道。实验中参考操作先改变待测电子态，再测量能量；数学上也可把这一过程等价地写成原始未知态上的旋转测量算符。扫描参考相位后，不同密度矩阵副对角线在相位傅里叶空间中被分离。

在有限能量窗口内，用整数 $\ell,n\in\mathbb Z$ 标记电子能量格基 $|\ell\rangle$。这些基矢无量纲且正交，未知电子态写成
\begin{equation*}
\hat\rho_e
=\sum_{\ell,n}\rho_{\ell n}|\ell\rangle\langle n|,
\qquad
\rho_{\ell n}\equiv\langle\ell|\hat\rho_e|n\rangle,
\end{equation*}
其中 $\rho_{\ell n}$ 无量纲，并满足 $\hat\rho_e\succeq0$、$\rho_{n\ell}=\rho_{\ell n}^*$ 与 $\sum_\ell\rho_{\ell\ell}=1$。对角元 $\rho_{\ell\ell}$ 是能量格占据概率，非对角元 $\rho_{\ell n}$ 则给出相隔 $|\ell-n|$ 个能量量子的相干幅度。

取一个已标定的平移协变参考幺正算符 $\hat U_R(\varphi_R)$。扫描参数 $\varphi_R\in[0,2\pi)$ 是无量纲相位；平移协变性意味着矩阵元只依赖末态与初态格点之差，因此定义无量纲参考振幅 $u_r$：
\begin{equation}
\boxed{
\langle m|\hat U_R(\varphi_R)|\ell\rangle
=u_{m-\ell}\,\ee^{\ii(m-\ell)\varphi_R}.}
\label{eq:reference-pinem-matrix}
\end{equation}
这里 $r=m-\ell\in\mathbb Z$ 是参考操作引起的净能量格位移。由于 $\hat U_R$ 幺正，这些系数不能任意选择，而必须满足
\begin{equation*}
\sum_m u_{m-\ell}u_{m-n}^*=\delta_{\ell n}.
\end{equation*}
实验上，$u_r$ 由已知参考场的耦合强度、相互作用长度和电子速度等标定参数决定；层析时它们被视为前向模型中已知的复数系数。

参考操作后理想测量末态能量格点 $m$。对应 POVM 元素为 $|m\rangle\langle m|$，Born 规则给出无量纲概率
\begin{equation}
\boxed{
P_m(\varphi_R)
=\sum_{\ell,n}
 u_{m-\ell}u_{m-n}^*\rho_{\ell n}
\ee^{\ii(n-\ell)\varphi_R}.}
\label{eq:electron_tomography_linear}
\end{equation}
\eqnrefs{eq:electron_tomography_linear} 已经是未知密度矩阵到可测能谱的线性前向模型。相位 $\varphi_R$ 的作用是给不同副对角线乘上不同 Fourier 谐波，因此对 $P_m(\varphi_R)$ 作相位 Fourier 分解可以逐条抽取副对角线。

定义第 $\Delta\ell\in\mathbb Z$ 阶 Fourier 系数
\begin{equation*}
\widetilde P_m^{(\Delta\ell)}
\equiv\frac{1}{2\pi}
\int_0^{2\pi}\dd\varphi_R\,
\ee^{-\ii (\Delta\ell)\varphi_R}P_m(\varphi_R).
\end{equation*}
$\widetilde P_m^{(\Delta\ell)}$ 仍无量纲。相位正交性只保留 $n-\ell=\Delta\ell$，于是
\begin{equation}
\boxed{
\widetilde P_m^{(\Delta\ell)}
=\sum_\ell \mathsf A_{m\ell}^{(\Delta\ell)}\,
\rho_{\ell,\ell+\Delta\ell},
\qquad
\mathsf A_{m\ell}^{(\Delta\ell)}
\equiv u_{m-\ell}u_{m-\ell-\Delta\ell}^*.}
\label{eq:electron_tomography_side_diagonal}
\end{equation}
$\mathsf A^{(\Delta\ell)}$ 是第 $\Delta\ell$ 条副对角线的\emph{层析设计矩阵}，其矩阵元无量纲。它由实验参考操作完全决定，而向量 $\rho_{\ell,\ell+\Delta\ell}$ 才是需要反演的未知量。

要判断某条副对角线是否真的可以由实验数据确定，不能只比较测量方程数目；必须检查 $\mathsf A^{(\Delta\ell)}$ 在目标参数子空间中的奇异值谱。写奇异值分解
\begin{equation}
\boxed{
\mathsf A^{(\Delta\ell)}
=\mathsf U_{\Delta\ell}\,
\operatorname{diag}(s_1^{(\Delta\ell)},s_2^{(\Delta\ell)},\ldots)
\mathsf V_{\Delta\ell}^\dagger,}
\label{eq:electron_tomography_svd}
\end{equation}
其中 $\mathsf U_{\Delta\ell},\mathsf V_{\Delta\ell}$ 为无量纲酉矩阵，$s_a^{(\Delta\ell)}\ge0$ 为无量纲奇异值。若 $s_a^{(\Delta\ell)}=0$，对应的右奇异向量位于测量映射零空间中：沿该密度矩阵方向改变未知态不会改变任何当前设置下的理想数据，因此该方向物理上不可辨识。若 $0<s_a^{(\Delta\ell)}\ll1$，伪逆会把数据噪声沿该方向放大约 $1/s_a^{(\Delta\ell)}$；这时改变参考耦合、增加传播设置或改变读出基可以改变设计矩阵并抬高最小有效奇异值，而单纯增加同一设置的重复计数只能降低统计噪声，不能移除零空间。

离散能量格层析适用于有限 PINEM/QPINEM 边带窗口。若实验通过二次电子光学连续扫描纵向相空间，则同一测量思想可写成 Wigner--Radon 层析。这里必须先把有量纲的纵向位置 $z$ 与动量 $p$ 化为一对正则无量纲变量。任选一个固定长度标尺 $z_0>0$，定义
\begin{equation}
\boxed{
X\equiv\frac{z}{z_0},
\qquad
P\equiv\frac{z_0p}{\hbar},
\qquad
[\hat X,\hat P]=\ii.}
\label{eq:dimensionless_longitudinal_phase_space}
\end{equation}
$z_0$ 的单位为米，用来选择相空间坐标尺度；可测结论对这一任意尺度的选择保持不变。有量纲 Wigner 函数 $W(z,p)$ 由\eqnrefs{eq:electron_wigner} 定义，并满足
\[
[W]=({\rm J\,s})^{-1},
\qquad
\int\dd z\dd p\,W(z,p)=1,
\]
因为 $\dd z\dd p$ 的单位正是作用量。对应的无量纲 Wigner 函数定义为
\begin{equation*}
W_{XP}(X,P)
\equiv\hbar\,W\!\left(z_0X,\frac{\hbar P}{z_0}\right),
\qquad
\int\dd X\dd P\,W_{XP}=1.
\end{equation*}

已标定的二次电子光学对 $(X,P)$ 实现线性辛旋转。用无量纲旋转角 $\theta$ 表示该设置，并定义被读出的旋转象限
\begin{equation*}
X_\theta=X\cos\theta+P\sin\theta.
\end{equation*}
若探测器理想测量 $X_\theta$，其结果记为无量纲变量 $x$，概率密度 $\mathcal P_\theta(x)$ 对 $x$ 归一化。Wigner 边缘性质给出
\begin{equation}
\boxed{
\mathcal P_\theta(x)
=\int\dd X\dd P\,
W_{XP}(X,P)
\delta[x-X\cos\theta-P\sin\theta].}
\label{eq:wigner_radon_transform}
\end{equation}
这就是 $W_{XP}$ 沿方向 $\theta$ 的 Radon 投影。

令 $s_X$ 为与无量纲读出变量 $x$ Fourier 共轭的无量纲 Fourier 变量，定义
\begin{equation*}
\widetilde{\mathcal P}_\theta(s_X)
\equiv\int\dd x\,\ee^{-\ii s_Xx}\mathcal P_\theta(x).
\end{equation*}
二维 Fourier 切片定理给出
\begin{equation}
\boxed{
\widetilde{\mathcal P}_\theta(s_X)
=\widetilde W(s_X\cos\theta,s_X\sin\theta),}
\label{eq:wigner_fourier_slice_theorem}
\end{equation}
其中 $\widetilde W(K_X,K_P)$ 是 $W_{XP}$ 的二维 Fourier 变换。于是扫描 $\theta$ 等价于在二维 Fourier 平面上采集不同径向切片。理想连续数据可由
\begin{equation}
\boxed{
W_{XP}(X,P)
=\frac{1}{(2\pi)^2}
\int_0^\pi\dd\theta
\int_{-\infty}^{\infty}\dd s_X\,|s_X|
\widetilde{\mathcal P}_\theta(s_X)
\ee^{\ii s_X(X\cos\theta+P\sin\theta)}}
\label{eq:wigner_filtered_backprojection}
\end{equation}
重建；其中 $|s_X|$ 来自二维 Fourier 平面改用有符号径向坐标后的雅可比行列式，因此高 $|s_X|$ 分量在反投影中被放大。

有限空间分辨率必须进入前向模型。令 $x_d$ 为探测器记录的无量纲读出，并定义条件线形 $r_X(x_d|x)$，满足
\begin{equation*}
r_X(x_d|x)\ge0,
\qquad
\int\dd x_d\,r_X(x_d|x)=1.
\end{equation*}
若响应平移不变，可写 $r_X(x_d|x)=r_X(x_d-x)$。观测投影为卷积，Fourier 空间中有
\begin{equation}
\boxed{
\widetilde{\mathcal P}^{\rm obs}_\theta(s_X)
=\widetilde r_X(s_X)\,
\widetilde{\mathcal P}_\theta(s_X).}
\label{eq:wigner_tomography_detector_transfer}
\end{equation}
当 $|\widetilde r_X(s_X)|$ 在高频趋近零时，形式反卷积 $1/\widetilde r_X(s_X)$ 会放大计数噪声；这与离散层析中小奇异值导致的噪声放大是同一个线性逆问题结构。

\medskip

自由漂移对纵向 Wigner 函数实施相空间剪切；不同漂移距离的投影等价于不同 Radon 角，从而允许层析。 建立无量纲相空间变量，推出 Radon 投影、Fourier 切片定理和带探测器传递函数的滤波反投影~\eqnrange{eq:dimensionless_longitudinal_phase_space}{eq:wigner_tomography_detector_transfer}；这里采用二次自由传播使映射线性；扫描角度覆盖所需范围；探测器传递函数在可反演频带内非零或用正则化处理。由 $X=z/z_0$、$P=z_0p/\hbar$ 与 $[\hat z,\hat p]=\ii\hbar$，
\[
[\hat X,\hat P]
=\frac1{z_0}\frac{z_0}{\hbar}[\hat z,\hat p]
=\ii.
\]
坐标变换的雅可比行列式为
\[
\left|\frac{\partial(z,p)}{\partial(X,P)}\right|
=z_0\frac{\hbar}{z_0}=\hbar,
\]
因此
\begin{align*}
\int\dd X\dd P\,W_{XP}(X,P)
&=\int\dd X\dd P\,
\hbar W\!\left(z_0X,\frac{\hbar P}{z_0}\right)\\
&=\int\dd z\dd p\,W(z,p)=1.
\end{align*}

辛旋转写成
\[
\begin{pmatrix}X_\theta\\P_\theta\end{pmatrix}
=\begin{pmatrix}
\cos\theta&\sin\theta\\
-\sin\theta&\cos\theta
\end{pmatrix}
\begin{pmatrix}X\\P\end{pmatrix}.
\]
旋转矩阵行列式为一，因此 $\dd X_\theta\dd P_\theta=\dd X\dd P$。测量 $X_\theta=x$ 而不读取 $P_\theta$ 时，
\begin{align*}
\mathcal P_\theta(x)
&=\int\dd P_\theta\,W_\theta(x,P_\theta)\\
&=\int\dd X_\theta\dd P_\theta\,
W_\theta(X_\theta,P_\theta)\delta(x-X_\theta)\\
&=\int\dd X\dd P\,W_{XP}(X,P)
\delta[x-X\cos\theta-P\sin\theta],
\end{align*}
得到\eqnrefs{eq:wigner_radon_transform}。对该式作一维 Fourier 变换：
\begin{align*}
\widetilde{\mathcal P}_\theta(s_X)
&=\int\dd x\,\ee^{-\ii s_Xx}\mathcal P_\theta(x)\\
&=\int\dd X\dd P\,W_{XP}(X,P)
\int\dd x\,\ee^{-\ii s_Xx}
\delta[x-X\cos\theta-P\sin\theta]\\
&=\int\dd X\dd P\,W_{XP}(X,P)
\ee^{-\ii s_X(X\cos\theta+P\sin\theta)}.
\end{align*}
若二维 Fourier 变换定义为
\[
\widetilde W(K_X,K_P)
=\int\dd X\dd P\,
\ee^{-\ii(K_XX+K_PP)}W_{XP}(X,P),
\]
逐项比较指数即得到
\[
K_X=s_X\cos\theta,
\qquad K_P=s_X\sin\theta,
\]
所以\eqnrefs{eq:wigner_fourier_slice_theorem} 成立。二维逆 Fourier 变换为
\[
W_{XP}(X,P)
=\frac1{(2\pi)^2}
\int\dd K_X\dd K_P\,
\widetilde W(K_X,K_P)
\ee^{\ii(K_XX+K_PP)}.
\]
使用带符号径向坐标
\[
K_X=s_X\cos\theta,
\qquad K_P=s_X\sin\theta,
\qquad 0\le\theta<\pi,
\qquad -\infty<s_X<\infty,
\]
雅可比行列式为
\begin{align*}
\det\frac{\partial(K_X,K_P)}{\partial(s_X,\theta)}
&=\det\begin{pmatrix}
\cos\theta&-s_X\sin\theta\\
\sin\theta&s_X\cos\theta
\end{pmatrix}\\
&=s_X.
\end{align*}
积分测度取绝对值，因此 $\dd K_X\dd K_P=|s_X|\dd s_X\dd\theta$。负 $s_X$ 已覆盖相反径向方向，所以 $0\le\theta<\pi$ 恰好覆盖 Fourier 平面一次。代入切片定理即得到\eqnrefs{eq:wigner_filtered_backprojection}。最后，若平移不变探测器线形满足
\[
\mathcal P_\theta^{\rm obs}(x_d)
=\int\dd x\,r_X(x_d-x)\mathcal P_\theta(x),
\]
则
\begin{align*}
\widetilde{\mathcal P}_\theta^{\rm obs}(s_X)
&=\int\dd x_d\,\ee^{-\ii s_Xx_d}
\int\dd x\,r_X(x_d-x)\mathcal P_\theta(x)\\
&=\int\dd x\,\mathcal P_\theta(x)\ee^{-\ii s_Xx}
\int\dd y\,r_X(y)\ee^{-\ii s_Xy}\\
&=\widetilde r_X(s_X)\widetilde{\mathcal P}_\theta(s_X),
\end{align*}
其中第二行令 $y=x_d-x$，得到\eqnrefs{eq:wigner_tomography_detector_transfer}。
\end{exbox}

\subsection{四个参考相位如何重建一个电子相干元}
\label{sec:electron-reference-tomography-worked}

上一节说明可以在能谱测量前加入参考操作，将相干信息变为概率变化。现在具体计算：未知电子态只占据两个通道时，扫描四个参考相位能确定哪些矩阵元。先保留一般的参考幺正算符 $U_R(\varphi_R)$，其中 $\varphi_R$ 是可控相位；将探测结果 $m$ 的效应算符记为 $E_m^{\rm det}$。测得概率为
\begin{equation}
 P_m(\varphi_R)=\operatorname{Tr}[\rho_e U_R^\dagger(\varphi_R)
 E_m^{\rm det}U_R(\varphi_R)].
 \label{eq:electron-reference-effect-rotation}
\end{equation}
$U_R$ 在实验中实际作用于电子；把它移到测量算符两侧，是同一过程的 Heisenberg 表述。\emph{层析}的任务是利用一组这样的已知效应算符，从概率反推出初始 $\rho_e$。自由电子的光学参考层析已有实验实现\cite{Priebe2017NatPhoton}；下面用一个有限支持模型算清最少需要读哪些数。

假设已独立验证未知态只占据 $|0\rangle,|1\rangle$，写
\begin{equation*}
 \rho_e=\begin{pmatrix}p&c_{01}\\c_{01}^*&1-p\end{pmatrix},\qquad
 0\le p\le1,\qquad |c_{01}|^2\le p(1-p).
\end{equation*}
按 $|0\rangle,|1\rangle$ 的顺序排列基矢，$p$ 是通道 $0$ 的占据概率，可以直接测得。厄米性要求两个非对角元互为共轭，所以只剩 $c_{01}=x_{01}+\ii y_{01}$ 的实部和虚部未知。由 $\det\rho_e=p(1-p)-|c_{01}|^2\ge0$，它们必须位于半径为 $\sqrt{p(1-p)}$ 的圆盘内；这个条件保证重建结果仍是合法量子态。

令参考为经典 PINEM 脉冲，幅度 $b_R=|\beta_R|$ 已知、相位 $\varphi_R$ 可扫描。在\emph{完整双边格}上有
\begin{equation*}
 \langle m|U_R|\ell\rangle
 =J_{m-\ell}(2b_R)\ee^{\ii(m-\ell)\varphi_R}.
\end{equation*}
未知态的两通道支持并不意味着参考操作也只有两通道：参考会把电子送到其他 $m$，不能把其 $2\times2$ 子块当作幺正算符。以下只统计末态 $m=0$ 的概率，同时把所有其他输出视为互补结果。

取 $u_0=J_0(2b_R)$、$u_1=J_1(2b_R)$，利用 $J_{-1}=-J_1$，由\eqnrefs{eq:electron_tomography_linear} 得
\begin{equation}
 \begin{aligned}
 P_0(\varphi_R)&=A_R-2u_0u_1\Re[c_{01}\ee^{\ii\varphi_R}],\\
 A_R&=u_0^2p+u_1^2(1-p).
 \end{aligned}
 \label{eq:electron-reference-two-channel-fringe}
\end{equation}
相位扫描将 $x_{01}$ 乘在余弦上、$y_{01}$ 乘在正弦上。四个相位 $0,\pi/2,\pi,3\pi/2$ 给出
\begin{equation}
 x_{01}=\frac{P_0(\pi)-P_0(0)}{4u_0u_1},\qquad
 y_{01}=\frac{P_0(\pi/2)-P_0(3\pi/2)}{4u_0u_1}.
 \label{eq:electron-reference-four-phase-inversion}
\end{equation}
这些式子要求 $u_0u_1\neq0$。若参考强度恰落在相关 Bessel 零点，即使计数无限多也无法用这个输出通道确定相干；接近零点时误差按 $1/|u_0u_1|$ 放大。这就是设计矩阵的小奇异值在一个具体实验中的含义。

\begin{lemma}[两通道参考干涉的正交相干读出]
设初态只支撑两个通道，参考操作到同一末态的两条振幅分别为实 $u_0$ 与 $-u_1e^{-\ii\varphi_R}$。末概率的相位依赖为 $A_R-2u_0u_1\operatorname{Re}(c_{01}e^{\ii\varphi_R})$。当 $u_0u_1\ne0$ 且参考幅及对角背景不随扫描变化时，四参考相位差分别线性恢复 $\operatorname{Re}c_{01}$ 和 $\operatorname{Im}c_{01}$，独立二项计数的方差按同一线性公式传播。
\end{lemma}
\begin{proof}
末态 $m=0$ 可以由初态 $0$ 不交换到达，振幅为 $u_0$；也可以由初态 $1$ 损失一阶到达，振幅为 $-u_1\ee^{-\ii\varphi_R}$。两条路径进入同一个末态，故交叉项是
\begin{equation*}
 u_0(-u_1\ee^{-\ii\varphi_R})^*\rho_{01}
 +(-u_1\ee^{-\ii\varphi_R})u_0\rho_{10}
 =-2u_0u_1\Re[c_{01}\ee^{\ii\varphi_R}].
\end{equation*}
加上各路径的对角概率得到\eqnrefs{eq:electron-reference-two-channel-fringe}。相差 $\pi$ 的两次读出具有相同背景 $A_R$、相反交叉项，作差后便分别留下 $x_{01}$ 和 $y_{01}$。二项方差通过同一个线性作差公式传播，给出\eqnrefs{eq:electron-reference-counting-uncertainty}。
\end{proof}


\begin{exbox}[四相位数据的解析重建与正性检查]
将四相位测得概率记为 $p_0,p_{\pi/2},p_\pi,p_{3\pi/2}$，已知输入占据 $p$ 和参考系数 $u_0,u_1$。求相干，并给出无需对角化的物理一致性条件。

先把条纹写成三个实系数：
\begin{equation*}
 P_0(\varphi_R)=A_R-C_R\cos\varphi_R+S_R\sin\varphi_R,
 \quad C_R=2u_0u_1x_{01},\quad S_R=2u_0u_1y_{01}.
\end{equation*}
四个设置依次读出 $A_R-C_R,A_R+S_R,A_R+C_R,A_R-S_R$。两两作差与作和即可解出
\begin{equation*}
 C_R=\frac{p_\pi-p_0}{2},\qquad
 S_R=\frac{p_{\pi/2}-p_{3\pi/2}}2,\qquad
 c_{01}=\frac{C_R+\ii S_R}{2u_0u_1}.
\end{equation*}
同一模型还要求
\begin{equation*}
 p_0+p_\pi=p_{\pi/2}+p_{3\pi/2}=2A_R,
 \qquad A_R=u_0^2p+u_1^2(1-p).
\end{equation*}
这是独立于未知相干的背景检查。正半定条件 $\det\rho_e\ge0$ 进一步给
\begin{equation*}
 C_R^2+S_R^2\le4u_0^2u_1^2p(1-p).
\end{equation*}
最后，由矩阵平方直接得到纯度
\begin{equation*}
 \operatorname{Tr}\rho_e^2=p^2+(1-p)^2
 +\frac{C_R^2+S_R^2}{2u_0^2u_1^2}.
\end{equation*}
等号边界对应纯态；$C_R=S_R=0$ 对应具有同样普通能谱的非相干混合。有限计数会使一致性关系近似成立，统计检验必须结合下述计数误差，不能把一次偏离直接当成模型失效。
\end{exbox}

有限计数时，设每个相位使用 $N$ 个独立电子，$\widehat P_0=N_0/N$。在理想完整记录条件下，$N_0$ 为二项计数，其方差为 $NP_0(1-P_0)$。四组电子独立，故
\begin{equation}
 \operatorname{Var}\widehat x_{01}
 =\frac{P_0(\pi)[1-P_0(\pi)]+P_0(0)[1-P_0(0)]}
 {16u_0^2u_1^2N},
 \label{eq:electron-reference-counting-uncertainty}
\end{equation}
$\widehat y_{01}$ 的式子把两相位换成 $\pi/2,3\pi/2$。用例题中的条纹系数，上式分子还可写成 $2A_R(1-A_R)-2C_R^2$；虚部分子把 $C_R$ 换成 $S_R$。若要求两个分量的标准差均不超过 $\varepsilon_c$，充分的每相位计数为
\begin{equation*}
 N\ge\frac{\max\{2A_R(1-A_R)-2C_R^2,\,
 2A_R(1-A_R)-2S_R^2\}}{16u_0^2u_1^2\varepsilon_c^2}.
\end{equation*}
总剂量为四组共 $4N$，另需用于确定 $p$ 和标定参考的计数。实际设计时可用 $P(1-P)\le1/4$ 给出与未知态无关的保守上界。

若存在能量误分配或漏检，应把 $E_m^{\rm det}$ 换成标定的效应算符重新构造设计矩阵。把只落在某个窗口内的计数重新归一化会引入依赖未知态的分母，一般不再服从上述线性公式。噪声导致线性反演略越出正性圆盘时，应在同一前向概率模型中作正半定约束的似然拟合，而不是把负本征值无说明地删掉。

\figref{fig:electron-reference-tomography} 显示同样的能谱如何对应不同参考条纹。
\begin{figure}[htbp]
 \centering
 \includegraphics[width=0.76\linewidth]{electron_reference_tomography.pdf}
 \caption{两通道电子态的参考读出。数值图取 $p=0.65$、$b_R=0.40$；实线使用 $c_{01}=0.20+0.15\ii$，虚线使用具有相同占据的非相干混合；四个标记点是解析反演所需的相位。误差棒表示每相位 $10^4$ 个电子的二项计数标准差，并非模拟测量数据。}
 \label{fig:electron-reference-tomography}
\end{figure}


